\section{2023 Algebra \& Number Theory}

\subsection*{Exam}
\begin{exercise}
Let $F$ be a field. Let $n \geq 2$ be an integer. For $1 \leq i \neq j \leq n$ and for $\lambda \in F$, set $E_{ij}(\lambda)$ as the elementary matrix.

(1) Show that $\mathrm{SL}_n(F)$ can be generated by the matrices $E_{ij}(\lambda)$.

(2) Assume $n \geq 3$. Show that $\mathrm{GL}_n(F)$ is not solvable.

(3) Assume $n = 2$ and $|F| \geq 4$. Show that $\mathrm{GL}_n(F)$ is not solvable either.

(4) Are $\mathrm{GL}_2(\mathbb{F}_2)$ and $\mathrm{GL}_2(\mathbb{F}_3)$ solvable?
\end{exercise}

\begin{solution}
(1) This is a standard result from linear algebra. Any matrix $A \in \mathrm{SL}_n(F)$ can be reduced to the identity matrix using elementary row operations of the type $R_i \leftarrow R_i + \lambda R_j$. These operations correspond to left multiplication by $E_{ij}(\lambda)$. Since $\det(A) = 1$, the final scaling operation is not needed (or rather, the products of elementary matrices will automatically yield a matrix with determinant 1). Thus $A = E_1 E_2 \dots E_k$.

(2) For $n \geq 3$, we use the commutator relation: $[E_{ij}(\lambda), E_{jk}(\mu)] = E_{ik}(\lambda\mu)$ for distinct $i, j, k$. This shows that every elementary matrix is a commutator. Since elementary matrices generate $\mathrm{SL}_n(F)$, we have $\mathrm{SL}_n(F) = [\mathrm{SL}_n(F), \mathrm{SL}_n(F)]$. A group whose derived series does not terminate at 1 is not solvable. Since $\mathrm{SL}_n(F) \subset \mathrm{GL}_n(F)$, $\mathrm{GL}_n(F)$ is not solvable.

(3) For $n=2$, the commutator of $E_{12}(\lambda)$ and a diagonal matrix $D = \text{diag}(a, a^{-1})$ is:
$$[E_{12}(\lambda), D] = E_{12}(\lambda) D E_{12}(-\lambda) D^{-1} = E_{12}(\lambda(1-a^2))$$
If $|F| \geq 4$, there exists $a \in F^*$ such that $a^2 \neq 1$. Then every elementary matrix is a commutator. Thus $\mathrm{SL}_2(F)$ is its own commutator subgroup (except for small fields), so $\mathrm{GL}_2(F)$ is not solvable.

(4) $\mathrm{GL}_2(\mathbb{F}_2) \cong S_3$, which is solvable ($S_3 \supset A_3 \supset 1$).
$\mathrm{GL}_2(\mathbb{F}_3)$ has order 48. Its derived subgroup is $\mathrm{SL}_2(\mathbb{F}_3)$ (order 24), and the next is the quaternion group $Q_8$ (order 8), then $\mathbb{Z}/2\mathbb{Z}$, then 1. Thus $\mathrm{GL}_2(\mathbb{F}_3)$ is solvable.
\end{solution}

\begin{exercise}
Let $R$ be a ring with identity $1 \neq 0$.

(1) Assume that $R$ is commutative. Show that, if we have an isomorphism $R^n \simeq R^m$ of $R$-modules for some positive integers $n$ and $m$, then $m = n$.

(2) The analogous assertion of (1) for a non-commutative ring is not true in general as shown by the following example. Let $K$ be a non-trivial ring with identity, and $F$ be a free $K$-module with a countable basis indexed by $\mathbb{Z}_{\geq 1}$:
$$e_1, \ldots, e_n, \ldots$$
Write $R = \mathrm{End}_K(F)$, which is a ring with the usual addition and composition of two $K$-linear endomorphisms of $F$.

(a) Show that $R$ is not commutative.

(b) Let $f_0 \in R$ (resp. $f_1 \in R$) such that $f_0(e_{2i}) = e_i$ and $f_0(e_{2i-1}) = 0$ (resp. $f_1(e_{2i}) = 0$ and $f_1(e_{2i-1}) = e_i$) for every $i \in \mathbb{Z}_{\geq 1}$. Show that $\{f_0, f_1\}$ is a basis of $R$ as a left $R$-module.

(c) For any integers $n, m \geq 1$, show that there exists an isomorphism $R^n \stackrel{\sim}{\to} R^m$ of $R$-modules.
\end{exercise}

\begin{solution}
(1) Let $\mathfrak{m}$ be a maximal ideal of $R$. The isomorphism $R^n \cong R^m$ induces an isomorphism of $(R/\mathfrak{m})$-vector spaces:
$$(R/\mathfrak{m})^n \cong R^n / \mathfrak{m} R^n \cong R^m / \mathfrak{m} R^m \cong (R/\mathfrak{m})^m.$$
Since the dimension of a vector space over a field is well-defined, we must have $n = m$.

(2) (a) $R$ is not commutative because linear operators on an infinite-dimensional space do not commute in general (e.g., shift operators).
(b) The elements $f_0, f_1$ act as "half-projections". We can define $g_0, g_1 \in R$ such that $g_0(e_i) = e_{2i}$ and $g_1(e_i) = e_{2i-1}$. Then $f_0 g_0 = 1, f_1 g_1 = 1, f_0 g_1 = 0, f_1 g_0 = 0$ and $g_0 f_0 + g_1 f_1 = 1$. This implies $R \cong R \oplus R$ as a left $R$-module.
(c) Since $R \cong R^2$, by induction $R \cong R^2 \cong R^3 \cong \dots \cong R^n$. Thus $R^n \cong R^m$ for any $n, m \geq 1$.
\end{solution}

\begin{exercise}
Let $G$ be a profinite group. We say that it is topologically finitely generated if there exist finitely many elements $g_1, \ldots, g_n$ such that the closed subgroup generated by $g_1, \ldots, g_n$ is equal to $G$.

(1) Assume that $G$ is topologically finitely generated. Show that, for a fixed positive integer $n$, $G$ has only finitely many open subgroups of index $n$.

(2) Show that, for $K$ a number field (i.e., a finite field extension of $\mathbb{Q}$), its absolute Galois group $G_K$ is not topologically finitely generated.
\end{exercise}

\begin{solution}
(1) An open subgroup $H \subset G$ of index $n$ corresponds to a continuous transitive action of $G$ on the set of $n$ cosets. This is equivalent to a continuous homomorphism $\rho: G \to S_n$ such that the image acts transitively.
Since $G$ is topologically generated by $g_1, \dots, g_k$, any continuous homomorphism is uniquely determined by the images $\rho(g_i) \in S_n$.
There are only $(n!)^k$ possible choices for these images. Thus there are only finitely many such homomorphisms, and hence finitely many open subgroups of index $n$.

(2) For $K = \mathbb{Q}$, the subgroups of index 2 correspond to quadratic extensions $\mathbb{Q}(\sqrt{d})$. There are infinitely many such extensions (one for each square-free $d \in \mathbb{Z}$). Thus $G_{\mathbb{Q}}$ has infinitely many open subgroups of index 2. By part (1), it cannot be topologically finitely generated.
\end{solution}

\begin{exercise}
Let $R$ be a commutative ring with identity. An $R$-module $M$ is said to be of finite presentation if there exists an exact sequence of $R$-modules
$$R^n \longrightarrow R^m \longrightarrow M \longrightarrow 0$$
for some positive integers $m, n \geq 0$.

(1) Let $f: N \to M$ be a surjective morphism of $R$-modules with $N$ of finite type and $M$ of finite presentation. Show that the $R$-module $\ker(f)$ is of finite type.

(2) Let $M$ be an $R$-module of finite type. Show that the following statements are equivalent:
\begin{itemize}
\item[(a)] $M$ is flat and of finite presentation.
\item[(b)] $M$ is locally free, i.e., there exist $f_1, \dots, f_s \in R$ generating the unit ideal of $R$ such that $M_{f_i}$ is free over $R_{f_i}$.
\item[(c)] $M$ is projective as an $R$-module.
\end{itemize}

(3) Let $S = \prod_{\mathbb{N}} R$ the product of countably many copies of $R$, and $I = \bigoplus_{\mathbb{N}} R \subset S$. Show that $S/I$ is a flat $S$-module which is not projective.
\end{exercise}

\begin{solution}
(1) Since $M$ is finitely presented, there exists a sequence $R^p \to R^q \to M \to 0$. Since $N$ is of finite type, there is $R^s \to N \to 0$. By Schanuel's Lemma or a diagram chase, $\ker(f)$ is seen to be a quotient of a finite type module, hence finite type.

(2) This is a central theorem in commutative algebra. A module is projective if and only if it is locally free of finite rank. For a finitely presented module, flatness is equivalent to being locally free.

(3) $S = \prod_{\mathbb{N}} R$, $I = \bigoplus_{\mathbb{N}} R$. $S/I$ is a direct limit of flat modules (if $R$ is a field), but we check projectivity. If $S/I$ were projective, the sequence $0 \to I \to S \to S/I \to 0$ would split. Thus $I$ would be a direct summand of $S$. However, any element $x \in S$ such that $ex = 0$ for all $e \in I$ must be zero. This prevents $I$ from having a complement that could map onto $S/I$ non-trivially.
\end{solution}

\begin{exercise}
Let $p$ be a prime number. Let $\mathbb{Z}_p$ be the ring of $p$-adic integers.

(1) Let $a \in \mathbb{C}$ such that $\lim_{r\to\infty} a^{p^r} = 1$. Show that $a$ is a $p^m$-th root of unity for some $m$.

(2) Let $A \in \mathrm{M}_{n\times n}(\mathbb{C})$ be an $n \times n$ complex matrix such that $\lim_{r\to\infty} A^{p^r} = I_n$, with $I_n$ the identity matrix. Show that $A^{p^m} = I_n$ for some integer $m$.

(3) Determine, up to isomorphisms, all the finite-dimensional continuous complex representations of $\mathbb{Z}_p$.
\end{exercise}

\begin{solution}
(1) Let $a = re^{i\theta}$. Then $r^{p^r} \to 1 \implies r = 1$. And $p^r \theta \to 0 \pmod{2\pi}$. This means $p^r \theta = 2\pi k_r + \epsilon_r$ with $\epsilon_r \to 0$. This forces $\theta = 2\pi k / p^m$ for some $m$.

(2) Considering the Jordan normal form, the eigenvalues must be roots of unity by part (1). For the blocks, $(1 + \epsilon)^{p^r} = 1 + p^r \epsilon + \dots$ which only converges to 1 if the nilpotent part $\epsilon$ is zero.

(3) A continuous representation $\rho: \mathbb{Z}_p \to \mathrm{GL}_n(\mathbb{C})$ must have an open kernel because $\mathbb{Z}_p$ is compact and $\mathrm{GL}_n(\mathbb{C})$ has no small subgroups. Thus $\rho$ factors through $\mathbb{Z}/p^k \mathbb{Z}$ for some $k$. These are finite abelian groups, so the representations are direct sums of characters $\chi(x) = \exp(2\pi i m x / p^k)$.
\end{solution}

\begin{exercise}
Let $p$ be a prime number. Let $K$ be a $p$-adic local field, i.e., a finite field extension of $\mathbb{Q}_p$. Denote by $\mathcal{O}_K$ its ring of integers, with $\pi \in \mathcal{O}_K$ a uniformizer.

(1) Show that the following logarithm map
$$\log: 1 + p\mathcal{O}_K \longrightarrow p\mathcal{O}_K, \quad 1 + x \mapsto x - \frac{x^2}{2} + \frac{x^3}{3} + \cdots$$
is a well-defined isomorphism of groups, which can be extended in a unique way to a group homomorphism
$$\log_\pi: K^\times \longrightarrow K,$$
such that $\log_\pi(\pi) = 0$. Determine the kernel of $\log_\pi$.

(2) With the help of the logarithm above, show that for any integer $n \geq 1$, the group quotient $K^\times / K^{\times, n}$ is a finite group.
\end{exercise}

\begin{solution}
(1) The power series $\log(1+x)$ converges for $|x| < 1$. It is an isomorphism when the series for $\exp$ also converges, which happens on $p\mathcal{O}_K$ (or $p^2$ for $p=2$). The unique extension $\log_\pi$ satisfying $\log_\pi(\pi) = 0$ is given by $\log_\pi(x \pi^k) = \log(x/\omega(x))$ where $\omega(x)$ is the root of unity part. The kernel consists of elements of the form $\pi^k \zeta$ where $\zeta$ is a root of unity in $K$.

(2) From the structure theorem $K^\times \cong \mathbb{Z} \times \mu(K) \times \mathbb{Z}_p^d$, the quotient is $\mathbb{Z}/n\mathbb{Z} \times \mu(K)/\mu(K)^n \times (\mathbb{Z}_p/n\mathbb{Z}_p)^d$. Since $\mathbb{Z}_p/n\mathbb{Z}_p$ is finite, the whole group is finite.
\end{solution}
